\chapter{Simplicial Complexes}
\url{https://team.inria.fr/geometrica/files/2015/11/1-simplicial-complexes.pdf}

\section{Geometric simplices}

\begin{defi}
    A \defini{$k$-simplex} $\sigma$ is the convex hull of $k+1$ points of $\R^d$ that are affinely independent. 
\end{defi}

\[
\sigma = \mathrm{conv}(p_0, \dots, p_k) = \ens[\bigg]{x \in \R^d, x = \sum_{i=0}^k \lambda_i \, p_i, \lambda_i \in \intervalleff{0}{1}, \sum_{i=0}^k \lambda_i = 1}
\]
$k = \dim(\mathrm{aff}(\sigma))$ is called the dimension of $\sigma$.

$1$-simple = line segment, $2$-simplex = triangle, $3$-simplex = tetrahedron

$V(\sigma)$ = set of vertices of a $k$-simplex $\sigma$

$\forall V' \subset V(\sigma)$, $\mathrm{conv}(V')$ is a face of $\sigma$.

\subsection{Geometric simplicial complexes}

A finite collection of simplices $K$ called the faces of $K$ such that
\begin{itemize}
    \item $\forall \sigma \in K$, $\sigma$ is a simplex
    \item $\sigma \in K$, $\tau \subset \sigma \implies \tau \in K$
    \item $\forall \sigma, \tau \in K$, either $\sigma \cap \tau = \varnothing$ or $\sigma \cap \tau$ is a common face of both
\end{itemize}

The dimension of a simplicial complex $K$ is the max dimension of its simplices.

The underlying space $\abs{K} \subset \R^d$ of $K$ is the union of the simplices of $K$. 

\subsection{Triangulation of a finite point set of \texorpdfstring{$\R^d$}{Rd}}

\begin{itemize}
    \item A simplicial $d$-complex $K$ is pure if every simplex in $K$ is the face of a $d$-simplex.
    \item A triangulation of a finite point set $P \in \R^d$ is a pure geometric simplicial complex $K$ s.t. $\mathrm{vert}(K) = P$ and $\abs{K} = \mathrm{conv}(P)$
\end{itemize}

\subsection{Triangulation of a polygonal domain of \texorpdfstring{$\R^2$}{R2}}

A triangulation of a polygonal domain $\Omega \in \R^2$ is a pure geometric simplicial complex $K$ s.t. $\mathrm{vert}(K) = \mathrm{vert}(\Omega)$ and $\abs{K} = \Omega$